
\documentclass{article}
\usepackage{colortbl}
\usepackage{makecell}
\usepackage{multirow}
\usepackage{supertabular}

\begin{document}

\newcounter{utterance}

\centering \large Interaction Transcript for game `eqbench', experiment `eq\_bench', episode 153 with reflexion{-}llama{-}3.1{-}8b{-}instruct.
\vspace{24pt}

{ \footnotesize  \setcounter{utterance}{1}
\setlength{\tabcolsep}{0pt}
\begin{supertabular}{c@{$\;$}|p{.15\linewidth}@{}p{.15\linewidth}p{.15\linewidth}p{.15\linewidth}p{.15\linewidth}p{.15\linewidth}}
    \# & \multicolumn{2}{c}{Player} && \multicolumn{2}{c}{Game Master} \\
    \hline

    \theutterance \stepcounter{utterance}  
    & & & \multicolumn{4}{p{0.6\linewidth}}{
        \cellcolor[rgb]{0.9,0.9,0.9}{
            \makecell[{{p{\linewidth}}}]{
                \texttt{\tiny{[P1$\langle$GM]}}
                \texttt{You have attempted this same task before and did not succeed. These are the plans you wrote after those attempts.} \\
\texttt{Trial \#1: New plan: I overestimated the intensity of Rev. Montague's emotions, particularly "Outraged" and "Indifferent". Upon re{-}reading the dialogue, I realize that Rev. Montague's tone remains calm and composed, even when faced with Dr. LeGuin's challenging arguments. He attempts to defend his views, but doesn't display any signs of outrage or indifference. In contrast, Rev. Montague is likely feeling challenged by Dr. LeGuin's criticisms, as evident from his defensive tone and the fact that he's trying to justify his views. I should have scored "Challenged" higher and "Outraged" lower. Next time, I will pay closer attention to the tone and language used by the characters to infer their emotional states.} \\
\\ 
\texttt{Trial \#2: New plan: I overestimated the intensity of Rev. Montague's emotions, particularly "Outraged" and "Indifferent". Upon re{-}reading the dialogue, I realize that Rev. Montague's tone remains calm and composed, even when faced with Dr. LeGuin's challenging arguments. However, I also realize that I should have considered the context of the dialogue more carefully. Rev. Montague is a religious figure, and his views are being challenged by a scientist. This could lead to a sense of being challenged, but also a sense of being defensive or even slightly hurt. I should have scored "Challenged" higher, but also considered the possibility of Rev. Montague feeling a bit defensive or hurt. Next time, I will pay closer attention to the context and the tone of the dialogue to infer the emotional states of the characters more accurately.} \\
\\ 
\\ 
\\ 
\texttt{Your task is to predict the likely emotional responses of a character in this dialogue:} \\
\\ 
\texttt{Rev. Montague: You see, Iris, faith is not a choice, but a divine gift. It's a beacon of light in a world often shrouded in darkness.} \\
\texttt{Dr. LeGuin: A beacon of light, Luther? Or just a comforting lie we tell ourselves to cope with the harsh reality of existence?} \\
\texttt{Rev. Montague: It's more than that. It's a guiding principle, a moral compass. It's what makes us human, Iris.} \\
\texttt{Dr. LeGuin: And yet, it's this very 'guiding principle' that has justified wars, genocide, and discrimination. How human is that, Luther?} \\
\texttt{Rev. Montague: Those are the failures of men, not faith. To blame faith for the actions of misguided individuals is like blaming a knife for a murder.} \\
\texttt{Dr. LeGuin: A knife doesn't claim to be a moral compass, Luther. And it doesn't demand blind obedience.} \\
\texttt{{[}End dialogue{]}} \\
\\ 
\texttt{At the end of this dialogue, Rev. Montague would feel...} \\
\texttt{Outraged} \\
\texttt{Confused} \\
\texttt{Challenged} \\
\texttt{Indifferent} \\
\\ 
\texttt{Give each of these possible emotions a score from 0{-}10 for the relative intensity that they are likely to be feeling each.} \\
\\ 
\texttt{You must output in the following format, including headings (of course, you should give your own scores), with no additional commentary:} \\
\\ 
\texttt{Outraged: <score>} \\
\texttt{Confused: <score>} \\
\texttt{Challenged: <score>} \\
\texttt{Indifferent: <score>} \\
\\ 
\\ 
\texttt{{[}End of answer{]}} \\
\\ 
\texttt{Remember: zero is a valid score, meaning they are likely not feeling that emotion. You must score at least one emotion > 0.} \\
\\ 
\texttt{Your answer:} \\
            }
        }
    }
     \\ \\

    \theutterance \stepcounter{utterance}  
    & \multicolumn{4}{p{0.6\linewidth}}{
        \cellcolor[rgb]{0.561,0.737,0.561}{
            \makecell[{{p{\linewidth}}}]{
                \texttt{\tiny{[P1$\rangle$GM]}}
                \texttt{Outraged: 2} \\
\texttt{Confused: 6} \\
\texttt{Challenged: 8} \\
\texttt{Indifferent: 0} \\
            }
        }
    }
    & & \\ \\

    \theutterance \stepcounter{utterance}  
    & & & \multicolumn{2}{p{0.3\linewidth}}{
        \cellcolor[rgb]{0.9,0.9,0.9}{
            \makecell[{{p{\linewidth}}}]{
                \texttt{\tiny{[GM$|$GM]}}
                \texttt{Outraged: 2} \\
\texttt{Confused: 6} \\
\texttt{Challenged: 8} \\
\texttt{Indifferent: 0} \\
            }
        }
    }
    & & \\ \\

    \theutterance \stepcounter{utterance}  
    & & & \multicolumn{2}{p{0.3\linewidth}}{
        \cellcolor[rgb]{0.9,0.9,0.9}{
            \makecell[{{p{\linewidth}}}]{
                \texttt{\tiny{[GM$|$GM]}}
                \texttt{Outraged: 3} \\
\texttt{Confused: 0} \\
\texttt{Challenged: 8} \\
\texttt{Indifferent: 0} \\
            }
        }
    }
    & & \\ \\

    \theutterance \stepcounter{utterance}  
    & & & \multicolumn{2}{p{0.3\linewidth}}{
        \cellcolor[rgb]{0.9,0.9,0.9}{
            \makecell[{{p{\linewidth}}}]{
                \texttt{\tiny{[GM$|$GM]}}
                \texttt{game\_result = LOSE} \\
            }
        }
    }
    & & \\ \\

    \theutterance \stepcounter{utterance}  
    & & \multicolumn{4}{p{0.6\linewidth}}{
        \cellcolor[rgb]{0.95,0.95,0.95}{
            \makecell[{{p{\linewidth}}}]{
                \texttt{\tiny{[P1$\leftrightarrow$P]}}
                \texttt{Post{-}episode Reflexion (not sent to GM):} \\
\texttt{New plan: I overestimated the intensity of Rev. Montague's emotions, particularly "Outraged" and "Indifferent". Upon re{-}reading the dialogue, I realize that Rev. Montague's tone remains calm and composed, even when faced with Dr. LeGuin's challenging arguments. However, I also realize that I should have considered the context of the dialogue more carefully. Rev. Montague is a religious figure, and his views are being challenged by a scientist. This could lead to a sense of being challenged, but also a sense of being defensive or even slightly hurt. I should have scored "Challenged" higher, but also considered the possibility of Rev. Montague feeling a bit defensive or hurt. Additionally, I should have paid closer attention to the language used by Dr. LeGuin, particularly her use of phrases like "comforting lie" and "blind obedience", which could be seen as provocative and potentially hurtful. Next time, I will pay closer attention to the context, tone, and language used in the dialogue to infer the emotional states of the characters more accurately. I will also consider the possibility of Rev. Montague feeling a mix of emotions, such as being challenged, defensive, and slightly hurt, rather than just one dominant emotion.} \\
            }
        }
    }
    & \\ \\

\end{supertabular}
}

\end{document}
